\chapter{Diagonal Extraction and De Pierro Majorisation}
\label{app:depierro_majorisation}
\label{subsubsec:depierro_gap}

This appendix expands the distinction between diagonal extraction and separable majorisation for the frozen prior curvature in Chapter~\ref{chapter:optimisation}.

Simply extracting a diagonal from the non-separable Hessian does not preserve the bound~\eqref{eq:jacobian_majoriser}. De Pierro's convexity construction instead gives a separable upper bound \cite{depierroModifiedExpectationMaximization1995a}. For each row of a chosen factorisation $\mathcal{B}^\top\mathrm{diag}\{\breve c_i\}\mathcal{B}$, with $\breve c_i\geq0$, the direct row inequality is
\begin{equation*}
    \left(\sum_j b_{ij}v_j\right)^2
    \leq
    \left(\sum_j |b_{ij}|\right)
    \sum_j |b_{ij}|v_j^2.
\end{equation*}
Writing $b_i=\sum_j|b_{ij}|$ gives the factorisation-specific comparison
\begin{equation}
    d_j^{\mathrm{DP}}
    =
    \sum_i |b_{ij}|\,b_i\,\breve c_i ,
    \qquad\text{compared with}\qquad
    d_j^{\mathrm{diag}}
    =
    \sum_i b_{ij}^2\,\breve c_i .
    \label{eq:depierro_vs_diag}
\end{equation}
For this chosen factorisation, $d_j^{\mathrm{DP}}\geq d_j^{\mathrm{diag}}$. An unmixed first-order difference row with exactly the two entries $\pm1/\Delta_\ell$ contributes a factor of $2$ relative to its extracted diagonal. Once $\Pi_{m,n}$ mixes difference channels, the row support and coefficients depend on the stencil, projector and factorisation, so no universal factor is implied.

The Gershgorin construction~\eqref{eq:diag_gersh_preconditioner} inflates along the \emph{modality} index only. It does not bound the spatial off-diagonal mass generated by the neighbour coupling in $\bm{D}^\top(\cdot)\bm{D}$, and therefore does not close the spatial majorisation gap. It is also one member of a family: for any $\pi>0$, the weighted arithmetic--geometric-mean bound gives $\operatorname{diag}(c_{11,r}+|c_{12,r}|/\pi,\;c_{22,r}+\pi|c_{12,r}|)\succeq C_r$, which is De Pierro's convexity device applied along the modality index rather than in space. The symmetric choice $\pi=1$ is used in Chapter~\ref{chapter:optimisation}; other choices redistribute the cross-term mass between \gls{pet} and \gls{spect}, and connect this construction to the modality-specific metric scaling in Equation~\eqref{eq:modality_precond_weight}.

An explicit De Pierro construction for a fixed factorisation would majorise the frozen quadratic prior surrogate, at the cost of a more conservative step. Because the implemented operand~\eqref{eq:diag_tight_preconditioner} is exactly $\operatorname{diag}(H_{\mathcal R}^{\mathrm{frozen}})$, the comparison in~\eqref{eq:depierro_vs_diag} applies to it directly: for the chosen factorisation, $d_j^{\mathrm{DP}}$ dominates the implemented Jacobi diagonal componentwise, and the difference is precisely the inflation that pays for dominating the discarded spatial couplings. No De Pierro model was implemented here.

\section{Separable Majorisation as Future Work}

The implemented Jacobi diagonal is the exact diagonal of the chosen frozen surrogate, so the scalar operand carries no approximation error relative to that quadratic's diagonal; the outstanding theoretical gap is majorisation. An extracted diagonal is not necessarily a majoriser of the spatially coupled quadratic, whereas a De Pierro separable diagonal constructed from an explicit factorisation as in Equation~\eqref{eq:depierro_vs_diag} provides a genuine separable upper bound and is therefore the appropriate route to an \glsentryshort{mm}-valid diagonal, although the bound may be more conservative.

For the coupled problem, a controlled comparison between the implemented Jacobi diagonal and such a De Pierro construction should also include grouped voxel-block variants that retain the cross-modal entries in Equation~\eqref{eq:block_tight_preconditioner_entries}, together with a proof of the corresponding Loewner bound for the grouped De Pierro construction. The evaluation should report preconditioner refresh time, peak temporary memory, majorisation slack relative to the frozen curvature on small grids, iteration and data-pass counts, and end-to-end reconstruction time under a common stopping criterion; no such benchmark was performed here.

